\chapter{Proxy applicatif — Squid et SquidGuard}

\section{Architecture}

\begin{figure}[H]
    \centering
    \fbox{\parbox{0.85\linewidth}{\centering
    \begin{verbatim}
Postes clients (V10, V20, V50, V60)
        ↓ HTTP/HTTPS (80/443)
   [pfSense Squid Transparent 3128/3129]
        ↓ intercept + SSL MITM
   [SquidGuard Filtrage URL]
        ↓
   [Internet]

V30 Informatique → accès direct (bypass proxy)
    \end{verbatim}
    \vspace{0.5cm}}}
    \caption{Architecture Proxy}
    \label{fig:proxy-arch}
\end{figure}

\section{Certificat CA pour SSL MITM}

\begin{lstlisting}
System > Certificate Manager > CAs > + Add

Descriptive name   : Proxy-CA-corpmsi
Method             : Create internal CA
Key type           : RSA
Key length         : 2048
Digest algorithm   : SHA256
Lifetime           : 3650 days
Common Name        : Proxy-CA-corpmsi
\end{lstlisting}

\textbf{Important :} Exporter le .crt et le déployer via GPO sur les VLANs contrôlés.

\section{Installation et configuration Squid}

\begin{lstlisting}
System > Package Manager > Available Packages
  → Rechercher "squid"
  → Installer
\end{lstlisting}

\begin{lstlisting}
Services > Squid Proxy Server > General

Enable Squid              : [X]
Proxy Port               : 3128
SSL Port                 : 3129 (MITM)
Keep Settings/Data       : [X]
Log Store Directory      : /var/squid/logs
\end{lstlisting}

\section{Mode Transparent}

\begin{lstlisting}
Transparent HTTP Proxy       : [X]
Transparent Proxy Interface  : DIRECTION, COMMERCIAL, RH, COMPTABILITE
Transparent HTTPS Proxy      : [X]
Bypass Proxy for Private     : [X]
\end{lstlisting}

\section{Inspection SSL}

\begin{lstlisting}
SSL Man In the Middle Filtering

Enable SSL filtering         : [X]
SSL Intercept Interfaces     : DIRECTION, COMMERCIAL, RH, COMPTABILITE
SSL Proxy Port              : 3129
SSL Certificate Authority   : Proxy-CA-corpmsi
Enable SSL Caching          : [X]
Certificate Adapt           : setValidAfter, setValidBefore
\end{lstlisting}

\subsection{Whitelist SSL (No Bump)}

\begin{lstlisting}
Target Categories > + Add

Name : Whitelist
Regular expression :
(^|\.)windowsupdate\.com$
(^|\.)update\.microsoft\.com$
(^|\.)download\.microsoft\.com$
(^|\.)login\.microsoftonline\.com$
(^|\.)office\.com$
(^|\.)office365\.com$
(^|\.)paypal\.com$
(^|\.)stripe\.com$
\end{lstlisting}

\section{Installation SquidGuard}

\begin{lstlisting}
System > Package Manager > Available Packages
  → Rechercher "squidguard"
  → Installer
\end{lstlisting}

\begin{lstlisting}
Services > SquidGuard Proxy Filter > General

Enable           : [X]
Enable log       : [X]
Blacklist        : [X]
Blacklist URL    : https://dsi.ut-capitole.fr/blacklists/download/blacklists.tar.gz
\end{lstlisting}

Télécharger (attendre 5–10 minutes).

\section{Catégories UT1 (après téléchargement)}

\begin{table}[H]
    \centering
    \begin{tabularx}{\textwidth}{|l|p{4cm}|}
        \hline
        \textbf{Catégorie} & \textbf{Description} \\
        \hline
        adult & Pornographie \\
        \hline
        gambling & Jeux d'argent \\
        \hline
        audio-video & Streaming (YouTube, Netflix) \\
        \hline
        malware & Virus, trojans \\
        \hline
        phishing & Hameçonnage \\
        \hline
        social\_networks & Facebook, Instagram, X \\
        \hline
        vpn & Services VPN anonymisants \\
        \hline
        download & Téléchargements (warez, torrents) \\
        \hline
    \end{tabularx}
    \caption{Catégories UT1 (partiel)}
    \label{tab:ut1-categories}
\end{table}

\section{ACLs par VLAN}

\subsection{ACL\_DIRECTION}

\begin{lstlisting}
Name        : ACL_DIRECTION
Client IP   : 192.168.10.0/24
Order       : 1

Configuration :
WL_Autorise (Whitelist)    → whitelist
adult / mixed_adult        → deny
gambling / arjel           → deny
audio-video / radio        → deny
social_networks            → allow
chat / webmail             → allow
download / warez           → deny
default                    → allow
\end{lstlisting}

\subsection{ACL\_RH}

\begin{lstlisting}
Name        : ACL_RH
Client IP   : 192.168.50.0/24
Order       : 3

Configuration :
WL_Autorise (Whitelist)    → whitelist
adult / mixed_adult        → deny
gambling / arjel           → deny
audio-video / radio        → deny
social_networks            → deny
chat / webmail             → deny
download / warez           → deny
default                    → allow
\end{lstlisting}

\subsection{ACL\_COMPTABILITE}

\begin{lstlisting}
Name           : ACL_COMPTABILITE
Client IP      : 192.168.60.0/24
Order          : 4
Default Action : deny

Configuration :
WL_Autorise (Whitelist)    → whitelist
(tous les autres)          → deny
default                    → deny
\end{lstlisting}

\section{NAT Port Forward}

Pour chaque VLAN contrôlé (DIRECTION, COMMERCIAL, RH, COMPTABILITE) :

\begin{lstlisting}
Firewall > NAT > Port Forward > + Add

Interface      : <VLAN>
Protocol       : TCP
Source         : <VLAN> net
Destination    : !(not) 192.168.0.0/8
Dest Port      : 80
Redirect IP    : 127.0.0.1
Redirect Port  : 3128
Description    : Redirect HTTP vers Squid
\end{lstlisting}

Répéter pour port 443 → 3129.

\section{Règles Firewall}

Pour chaque interface contrôlée :

\begin{lstlisting}
Firewall > Rules > <VLAN> > + Add

Action     : PASS
Protocol   : TCP
Source     : <VLAN> net
Destination: 127.0.0.1
Port dest  : 3128, 3129
Description: Autoriser Squid

Action     : BLOCK
Protocol   : TCP
Source     : <VLAN> net
Destination: !(not) 192.168.0.0/8
Port dest  : 80, 443
Description: Bloquer bypass proxy
\end{lstlisting}

\section{Page de blocage}

\begin{lstlisting}[language=bash]
cat > /usr/local/www/blocked.html << 'HTML'
<!DOCTYPE html>
<html>
<head>
    <title>Accès bloqué</title>
    <style>
        body { font-family: Arial; text-align: center; padding: 50px; }
        .box { background: white; padding: 40px; max-width: 500px;
               margin: auto; border-radius: 8px; box-shadow: 0 2px 10px rgba(0,0,0,0.1); }
        h1 { color: #d9534f; }
    </style>
</head>
<body>
    <div class="box">
        <h1>Accès refusé</h1>
        <p>Ce site est bloqué par la politique de sécurité de <strong>corpmsi</strong>.</p>
        <p>Contactez le service informatique pour une demande d'exception.</p>
    </div>
</body>
</html>
HTML
\end{lstlisting}

\section{Récapitulatif des politiques par VLAN}

\begin{table}[H]
    \centering
    \begin{tabularx}{\textwidth}{|l|c|c|c|c|}
        \hline
        \textbf{Catégorie} & \textbf{V10 Dir} & \textbf{V20 Com} & \textbf{V50 RH} & \textbf{V60 Compt} \\
        \hline
        Social networks & allow & allow & deny & deny \\
        \hline
        Webmail & allow & allow & deny & deny \\
        \hline
        Chat & allow & allow & deny & deny \\
        \hline
        Audio-video & deny & deny & deny & deny \\
        \hline
        Gambling & deny & deny & deny & deny \\
        \hline
        Adult & deny & deny & deny & deny \\
        \hline
        Download & deny & deny & deny & deny \\
        \hline
        Default & allow & allow & allow & deny \\
        \hline
    \end{tabularx}
    \caption{Politiques proxy par VLAN}
    \label{tab:proxy-policies}
\end{table}

\section{Checklist}

\begin{itemize}
    \item[\checkmark] Créer CA Proxy-CA-corpmsi
    \item[\checkmark] Exporter .crt et déployer via GPO
    \item[\checkmark] Installer Squid + SquidGuard
    \item[\checkmark] Configurer mode transparent
    \item[\checkmark] Activer SSL MITM
    \item[\checkmark] Télécharger blacklists UT1
    \item[\checkmark] Créer ACLs par VLAN
    \item[\checkmark] Configurer NAT port forward
    \item[\checkmark] Ajouter règles firewall
    \item[\checkmark] Tester depuis chaque VLAN
\end{itemize}
